\section{An identifiability obstruction without target information}
\label{sec:p3-source-identifiability}
The observed lack of repeatable source-only benefits should
not be interpreted as a universal theorem that no EEG
domain-generalisation method can succeed. A more precise
mathematical obstruction explains why no \emph{unconditional}
guarantee can be established from source observations
alone.

\paragraph{Proposition.}
Let two fixed source-fitted binary linear policies
have nonparallel normal vectors $w_a,w_b$.
Without restrictions on the distribution of a
previously unseen target domain, the sign of
\[
 \mathrm{AUC}_{T}(w_a)-\mathrm{AUC}_{T}(w_b)
\]
is not identifiable from their identical source-only
training observations. Specifically, there exist
two target joint distributions $T_+$ and $T_-$
for which the source observations and both
fitted policies are unchanged but the pairwise
AUC preference reverses.

\paragraph{Construction.}
Since $w_a$ and $w_b$ are linearly independent
up to scaling, there exists a direction $v$
with $w_a^\top v>0$ and $w_b^\top v<0$.
Take target world $T_+$ with positive examples
at $v$ and negative examples at zero.
For target world $T_-$, exchange their
class labels. The linear ranking preference
changes from
$\mathrm{AUC}_{T_+}(w_a)=1$,
$\mathrm{AUC}_{T_+}(w_b)=0$ to
$\mathrm{AUC}_{T_-}(w_a)=0$,
$\mathrm{AUC}_{T_-}(w_b)=1$.
Arbitrarily small continuous perturbations
preserve the strict inequalities.
No source-only procedure observing the
same source law in both target worlds
can correctly predict both signs.
This is an elementary construction and
not a claim of a new no-free-lunch theorem.

The construction is reproduced by a
four-pair synthetic numerical regression.
It is a falsification of a proposed
\emph{distribution-free target sign
certificate}, not an experimental EEG result
or a clinical forecast.

\paragraph{What can make forecasting meaningful?}
Additional structural assumptions could include
restricted shifts of nuisance locations,
invariant conditional class contrast,
a bounded divergence between source and target
feature distributions, calibration EEG from the
target, or a validated source-to-target relationship.
Each assumption must be stated, measured and
tested on genuinely new acquisitions.
In particular, the group-centred discriminant
operator's invariance under fixed-coordinate
source-group translations does not imply
invariance under a newly refitted Riemannian
reference nor unchanged target label relations.
The failed Lee2019 positive-benefit promotion
is consistent with this limitation but
is not a general impossibility proof.

\section{How to turn the current negative case series into a falsifiable study}
The current four source-only algorithms were not
one predeclared candidate family tested
on a common external data matrix.
A suitable future primary endpoint would be
the \emph{source-only selected regret},
defined on truly new target groups $g$ by
\[
 R_g=\max_{m\in\mathcal M}
        \mathrm{AUC}_g(m)-
        \mathrm{AUC}_g(\widehat m_{\rm source}),
\]
where $\mathcal M$ and the source-only
selection rule are frozen before any target
EEG is read. This regret uses target labels
\emph{only for final outcome analysis},
not for source model selection.
A second endpoint is the rate of
$\mathrm{AUC}_g(\widehat m_{\rm source})-
\mathrm{AUC}_g(m_0)<-0.02$ against an
explicit original source-macro baseline.
The same candidates, channel selections
and source-locked participant units must
be evaluated on all external datasets.
Until this experiment exists, we report
the four earlier case studies as negative
research evidence rather than a replicated
performance theorem.
